\clearpage\pagestyle{matSimple}\phantomsection\label{guide}\pdfbookmark[0]{Guide de l’adulte}{bm-guide}
\matPageTitre{Guide de l’adulte}{Objectifs, manipulations, mots à faire dire et réponses}
{\fontsize{9.5}{13.5}\selectfont Pour chaque activité : l’objectif tiré du programme, ce que l’on peut faire avec de vrais objets avant la fiche, les mots que l’enfant est invité à dire, puis la réponse attendue. Les images en vert montrent une réponse correcte ; d’autres réponses peuvent être justes lorsque la consigne le permet.\par}
\par\needspace{62mm}\vspace{6mm}{\otSemi\fontsize{9}{11}\selectfont\color{matAccent}\otCaps[8]{Partie 1 · Plus, moins, autant}}\par\vspace{1mm}
\guideEntree{act-ms-1}{1}{Cinq partout}{Reconnaître des collections de cinq objets de tailles et de dispositions différentes.}{Faire des tas de cinq objets très différents (cinq perles, cinq livres, cinq ballons de tailles variées) et vérifier en comptant.}{« Les ballons sont gros ou petits, mais il y en a cinq. »}{Les boîtes des étoiles, des ballons, des poissons et des jetons rouges : quatre boîtes. Les pommes sont quatre, les jetons violets six.}{ms-01}
\guideEntree{act-ms-2}{2}{Le plus de billes}{Comparer globalement deux collections dont l’une a au moins deux fois plus d’objets que l’autre.}{Remplir deux pots transparents, l’un avec beaucoup de perles, l’autre avec peu ; demander sans compter où il y en a le plus.}{« Dans ce bocal, il y a plus de billes que dans l’autre. Dans l’autre, il y en a moins. »}{Le bocal de droite (dix-huit billes contre sept). On ne demande pas de compter : la différence se voit.}{ms-02}
\guideEntree{act-ms-3}{3}{Autant de jetons ?}{Comprendre qu’une quantité ne dépend pas de la place des objets : une rangée plus longue n’a pas forcément plus d’objets.}{Aligner cinq jetons rouges bien espacés et cinq bleus serrés ; demander où il y en a le plus, puis vérifier en les appariant.}{« La rangée rouge est plus longue, mais il y a autant de jetons : cinq et cinq. »}{En haut : autant (cinq et cinq). En bas : il y a plus de jetons bleus (six) que de rouges (quatre), même si la rangée bleue est plus courte.}{ms-03}
\guideEntree{act-ms-4}{4}{Un verre par assiette}{Dénombrer une collection pour en constituer une autre de même quantité, sans pouvoir les rapprocher.}{Mettre des assiettes sur une table et des verres sur une autre, loin : l’enfant doit rapporter juste assez de verres en un seul voyage.}{« Il y a cinq assiettes. Je prends cinq verres : un, deux, trois, quatre, cinq. »}{Cinq verres coloriés, n’importe lesquels : il y a cinq assiettes.}{ms-04}
\guideEntree{act-ms-5}{5}{Voitures ou bateaux ?}{Comparer deux collections en dénombrant chacune d’elles.}{Poser quatre voitures sur une table et six bateaux sur une autre table, loin : comparer en comptant chaque groupe.}{« Il y a quatre voitures et six bateaux. Il y a plus de bateaux que de voitures. »}{Les bateaux : six bateaux contre quatre voitures, même si les voitures prennent plus de place.}{ms-05}
\par\needspace{62mm}\vspace{6mm}{\otSemi\fontsize{9}{11}\selectfont\color{matAccent}\otCaps[8]{Partie 2 · Compter jusqu’à six}}\par\vspace{1mm}
\guideEntree{act-ms-6}{6}{Le dé et les boîtes}{Dénombrer une collection jusqu’à six et l’associer à la constellation du dé.}{Lancer un grand dé, compter les points, puis prendre autant d’objets dans une réserve.}{« Il y a quatre étoiles : je cherche le dé à quatre points. »}{Quatre étoiles avec le dé quatre, six jetons avec le dé six, cinq fleurs avec le dé cinq, trois pommes avec le dé trois.}{ms-06}
\guideEntree{act-ms-7}{7}{Doigts et chiffres}{Associer une quantité montrée avec les doigts et son écriture chiffrée.}{Montrer un nombre avec ses doigts ; l’enfant le dit, puis le retrouve parmi des cartes chiffrées.}{« Cinq doigts sur une main et un sur l’autre : six. Je relie au 6. »}{Six doigts avec 6, quatre avec 4, cinq avec 5, trois avec 3. Toute façon de montrer un nombre avec ses doigts est juste.}{ms-07}
\guideEntree{act-ms-8}{8}{Combien de points ?}{Dénombrer une collection mal rangée en pointant chaque élément une seule fois.}{Compter des jetons posés en vrac en les déplaçant un à un d’un côté à l’autre de la table.}{« Je touche chaque point une fois : un, deux, trois, quatre, cinq. Il y en a cinq. »}{Cinq, puis trois, puis six, puis quatre.}{ms-08}
\guideEntree{act-ms-9}{9}{Autant que le dé}{Constituer une collection de même quantité qu’une constellation donnée.}{Jeu de marchande : lancer le dé et acheter autant de fruits (jetons) qu’il y a de points.}{« Quatre points : je dessine quatre pommes. Je compte pour vérifier. »}{Quatre, six et cinq pommes. Des ronds suffisent pour représenter les pommes.}{ms-09}
\guideEntree{act-ms-10}{10}{Un ballon de plus}{Réaliser une collection qui a un objet de plus et nommer le nombre avant et après l’ajout.}{Faire une tour de trois cubes, en ajouter un et dire « trois, et encore un, quatre ».}{« Il y avait trois ballons, j’en dessine un de plus : il y en a quatre. »}{Trois et un, quatre ; cinq et un, six ; quatre et un, cinq. Quand on ajoute un, on dit le nombre suivant.}{ms-10}
\guideEntree{act-ms-11}{11}{Autant que le chiffre}{Passer de l’écriture chiffrée d’un nombre à la quantité.}{Tirer une carte chiffrée, la lire avec l’adulte, puis prendre autant de cubes.}{« C’est un quatre : je colorie quatre cases. »}{Quatre, six, deux et cinq cases coloriées, à partir de la gauche ou dans n’importe quel ordre.}{ms-11}
\guideEntree{act-ms-12}{12}{Le papillon butine}{Parcourir une collection en passant une et une seule fois par chaque élément.}{Poser un jeton dans chaque boîte d’allumettes fermée, sans en oublier ni en mettre deux, en déplaçant les boîtes déjà servies.}{« Celle-là, je l’ai déjà visitée. Il me reste celle-ci. »}{Il y a dix fleurs. Beaucoup de chemins sont justes : chacun doit toucher toutes les fleurs, une seule fois.}{ms-12}
\par\needspace{62mm}\vspace{6mm}{\otSemi\fontsize{9}{11}\selectfont\color{matAccent}\otCaps[8]{Partie 3 · Composer et écrire}}\par\vspace{1mm}
\guideEntree{act-ms-13}{13}{Les doigts des deux mains}{Composer des nombres jusqu’à six avec les doigts des deux mains.}{Montrer trois doigts sur une main et deux sur l’autre : combien en tout ? Recommencer avec d’autres façons.}{« Trois et deux, ça fait cinq. Quatre et deux, ça fait six. »}{Cinq, six, cinq et six. Remarquer que trois et deux font cinq, comme quatre et un.}{ms-13}
\guideEntree{act-ms-14}{14}{Cinq perles, deux couleurs}{Décomposer cinq de plusieurs façons : un et quatre, deux et trois…}{Faire des tours de cinq cubes avec deux couleurs et comparer toutes les tours obtenues.}{« Mon collier a deux perles rouges et trois bleues : cinq, c’est deux et trois. »}{Par exemple : un rouge et quatre bleus, deux et trois, trois et deux, quatre et un. L’ordre des perles peut varier.}{ms-14}
\guideEntree{act-ms-15}{15}{J’écris 1, 2, 3}{Écrire en chiffres les nombres de un à six.}{Tracer le chiffre en grand dans l’air, sur le dos de l’enfant, dans le sable ou sur l’ardoise, en partant du haut.}{« Je pars du haut, comme le modèle bleu. »}{Les chiffres sont tracés en partant du haut, sans lever le crayon (sauf pour le 4 et le 5). La forme compte plus que la taille.}{ms-15}
\guideEntree{act-ms-16}{16}{J’écris 4, 5, 6}{Écrire en chiffres les nombres de un à six.}{Tracer le chiffre en grand dans l’air, sur le dos de l’enfant, dans le sable ou sur l’ardoise, en partant du haut.}{« Je pars du haut, comme le modèle bleu. »}{Les chiffres sont tracés en partant du haut, sans lever le crayon (sauf pour le 4 et le 5). La forme compte plus que la taille.}{ms-16}
\guideEntree{act-ms-17}{17}{J’écris combien}{Représenter une quantité par son écriture chiffrée.}{Jeu du facteur : compter des jetons et écrire le nombre sur une étiquette pour un camarade qui doit en apporter autant.}{« Il y a deux canards : j’écris le chiffre deux. »}{Les chiffres 2, 5, 3 et 6.}{ms-17}
\guideEntree{act-ms-18}{18}{Les pierres de la rivière}{Réciter la comptine numérique de un à douze, de façon ordonnée, un mot par pierre.}{Réciter la comptine en sautant de cerceau en cerceau, en tapant dans les mains, en montant un escalier.}{« Un, deux, trois… dix, onze, douze ! Le lapin a traversé. »}{Douze pierres : un mot-nombre par saut, sans en oublier. La dernière pierre, près de la carotte, est la douzième.}{ms-18}
\par\needspace{62mm}\vspace{6mm}{\otSemi\fontsize{9}{11}\selectfont\color{matAccent}\otCaps[8]{Partie 4 · Rangs et cases}}\par\vspace{1mm}
\guideEntree{act-ms-19}{19}{En partant de la mare}{Repérer le rang d’un objet dans une file en comptant, à partir d’un point de départ donné.}{Faire une file d’enfants ou de peluches : « Qui est le premier ? Le deuxième ? Va te mettre en cinquième place. »}{« Je pars de la mare : premier le canard, deuxième le lapin… »}{Le lapin est le deuxième ; la souris est la cinquième. En partant de l’arbre, les rangs changent : la souris serait la deuxième.}{ms-19}
\guideEntree{act-ms-20}{20}{Le train des rangs}{Repérer un rang jusqu’au sixième dans une suite ordonnée.}{Aligner six chaises derrière une « locomotive » ; demander aux enfants de s’asseoir sur la première, la troisième, la sixième.}{« Le premier wagon est collé à la locomotive. Un, deux, trois : voilà le troisième. »}{Le wagon juste derrière la locomotive en rouge, le troisième en bleu, le dernier (sixième) en vert.}{ms-20}
\guideEntree{act-ms-21}{21}{La perle rouge}{Repérer le rang d’un élément en comptant à partir d’un point de départ.}{Enfiler cinq perles bleues et une rouge ; un enfant décrit la place de la rouge à un camarade qui doit refaire le même collier.}{« Je pars du nœud : un, deux, trois. La perle rouge est la troisième. »}{Le deuxième fil. Sur les autres, la perle rouge est la quatrième, la première et la cinquième.}{ms-21}
\guideEntree{act-ms-22}{22}{La bande des nombres}{Compléter une bande numérique lacunaire jusqu’à six.}{Construire avec la classe une bande de boîtes : une boîte par nombre, de un à six, avec la quantité et le chiffre.}{« Après un, il y a deux. Avant six, il y a cinq. »}{Première bande : 2 et 5. Deuxième : 3, 4 et 6. Troisième : 1, 4 et 5.}{ms-22}
\guideEntree{act-ms-23}{23}{Chaque carte à sa place}{Placer des représentations des nombres (dé, doigts, points) sur la bande numérique.}{Distribuer des cartes (dés, doigts, chiffres) ; chaque enfant pose la sienne sur la bonne case d’une grande bande au sol.}{« Quatre points : je vais à la case du 4. »}{Le dé quatre va au 4, les deux doigts au 2, la carte à six points au 6, le dé trois au 3, la main ouverte au 5, le point seul au 1.}{ms-23}
\guideEntree{act-ms-24}{24}{Le jeu de l’oie}{Trouver la position d’arrivée après un avancement d’une ou deux cases.}{Jouer au jeu de l’oie avec un dé à un ou deux points ; avancer le pion en comptant chaque case à voix haute.}{« Je suis ici. Le dé fait deux : un, deux. J’arrive là. »}{Le pion arrive deux cases plus loin, puis une case plus loin, puis deux cases plus loin. On ne compte pas la case de départ.}{ms-24}
\guideEntree{act-ms-25}{25}{On recule !}{Trouver la position d’arrivée après un recul d’une ou deux cases.}{Au jeu de l’oie, ajouter des cases « recule d’une case » ou « recule de deux cases ».}{« Je recule de deux : un, deux. Je suis revenu là. »}{Le pion recule d’une case, puis de deux, puis de deux.}{ms-25}
\par\needspace{62mm}\vspace{6mm}{\otSemi\fontsize{9}{11}\selectfont\color{matAccent}\otCaps[8]{Partie 5 · Petits problèmes}}\par\vspace{1mm}
\guideEntree{act-ms-26}{26}{Les fraises de Lilou}{Résoudre un problème de retrait en représentant l’action.}{Jouer la scène avec cinq jetons : en enlever deux, puis compter ce qui reste.}{« Il y avait cinq fraises, elle en a mangé deux : il en reste trois. Cinq, c’est deux et trois. »}{Deux fraises barrées (n’importe lesquelles) ; il en reste trois : le dé à trois points.}{ms-26}
\guideEntree{act-ms-27}{27}{Le chapeau du magicien}{Trouver une partie quand on connaît le tout et l’autre partie.}{Compter six jetons, puis en cacher sous un gobelet pendant que l’enfant ferme les yeux : combien sont cachés ?}{« Il y en a six en tout. J’en vois quatre : cinq, six… il en manque deux. »}{Deux lapins sont cachés : six, c’est quatre et deux.}{ms-27}
\guideEntree{act-ms-28}{28}{Les billes dans la boîte}{Résoudre un problème d’ajout : trouver le tout.}{Mettre trois billes dans une boîte opaque, en ajouter deux devant l’enfant et demander combien il y en a sans ouvrir.}{« Trois dans ma tête, et encore deux : quatre, cinq. Il y en a cinq. »}{Cinq billes : on peut surcompter à partir de trois (« quatre, cinq ») ou compter sur ses doigts.}{ms-28}
\guideEntree{act-ms-29}{29}{Sous la main}{Rechercher une partie dans un problème de parties-tout.}{Jeu du cache : poser cinq jetons, en cacher sous sa main, l’enfant trouve combien sont cachés, puis on vérifie.}{« Il y en a cinq en tout. J’en vois trois : il en manque deux. »}{Deux jetons cachés, puis deux, puis trois.}{ms-29}
\guideEntree{act-ms-30}{30}{Les boîtes de crayons}{Rechercher le tout dans un problème de groupements.}{Montrer trois boîtes opaques contenant chacune deux crayons, l’une après l’autre ; demander combien il y a de crayons en tout.}{« Deux, et deux, et encore deux : six crayons. »}{Six crayons : le dé à six points.}{ms-30}
\guideEntree{act-ms-31}{31}{Le partage des gâteaux}{Trouver la valeur d’une part dans un partage équitable, en distribuant un à un.}{Distribuer six jetons à deux poupées : « un pour toi, un pour toi », jusqu’à ce qu’il n’en reste plus.}{« Un pour toi, un pour toi… Chaque ours a trois gâteaux. »}{Trois gâteaux dans chaque assiette.}{ms-31}
\guideEntree{act-ms-32}{32}{Une part pour chaque lapin}{Partager équitablement une collection en distribuant un à un.}{Distribuer six cubes à trois peluches, un par un, en faisant le tour, jusqu’à ce qu’il n’en reste plus.}{« Un pour chacun, puis encore un pour chacun : chaque lapin a deux carottes. »}{Deux carottes pour chaque lapin.}{ms-32}
\par\needspace{62mm}\vspace{6mm}{\otSemi\fontsize{9}{11}\selectfont\color{matAccent}\otCaps[8]{Partie 6 · Formes et solides}}\par\vspace{1mm}
\guideEntree{act-ms-33}{33}{Triangles, carrés, disques}{Reconnaître et classer triangles, carrés et disques, quelles que soient leur taille et leur orientation.}{Trier des formes en carton de toutes tailles et tournées dans tous les sens ; les toucher les yeux fermés.}{« Ce carré est penché, mais c’est toujours un carré : quatre côtés pareils. »}{Quatre triangles en bleu, quatre carrés en rouge, quatre disques en jaune. Les carrés penchés sont bien des carrés.}{ms-33}
\guideEntree{act-ms-34}{34}{L’intrus}{Reconnaître une forme plane malgré sa couleur, sa taille et son orientation.}{Montrer quatre carrés tournés et un triangle ; demander lequel « n’est pas de la même famille ».}{« Ce sont tous des carrés, même tournés. Celui-là est un triangle : c’est l’intrus. »}{Le triangle de la première ligne, le carré de la deuxième, le carré de la troisième, le disque de la quatrième.}{ms-34}
\guideEntree{act-ms-35}{35}{Les empreintes du cube}{Produire les empreintes d’un solide et les reconnaître.}{Appuyer chaque face d’un cube dans la pâte à modeler ou la peinture : toutes les empreintes sont des carrés pareils.}{« Toutes les faces du cube laissent un carré. Le cube ne laisse pas de disque. »}{Les trois carrés, même tournés. Le rectangle, le disque et le triangle ne viennent pas d’un cube.}{ms-35}
\guideEntree{act-ms-36}{36}{Les empreintes de la pyramide}{Produire les empreintes d’un solide et retrouver le solide à partir d’une empreinte.}{Faire des empreintes avec une pyramide à base carrée : le dessous laisse un carré, les côtés laissent des triangles.}{« La pyramide laisse un carré et des triangles. Le cube aussi laisse un carré : une même empreinte peut venir de deux solides. »}{Le carré et les trois triangles. Le disque (empreinte d’un cylindre) et le rectangle ne viennent pas de cette pyramide.}{ms-36}
\guideEntree{act-ms-37}{37}{La fusée}{Reproduire un assemblage de cinq formes à partir d’un modèle.}{Reproduire la fusée avec cinq formes en carton posées sur une feuille, puis la coller.}{« En haut, le triangle rouge. Dessous, le carré bleu, puis le carré jaune. Les ailes sont vertes. »}{Pointe rouge, carré du haut bleu, carré du bas jaune, ailerons verts.}{ms-37}
\guideEntree{act-ms-38}{38}{Le carré en cinq pièces}{Reproduire un pavage de cinq pièces à partir d’un modèle.}{Découper un carré en cinq pièces de carton ; mélanger les pièces et reconstituer le carré en regardant le modèle.}{« Le grand triangle bleu est en bas à gauche. Le carré rouge est en haut à droite. »}{Chaque pièce a la couleur de la pièce placée au même endroit sur le modèle.}{ms-38}
\par\needspace{62mm}\vspace{6mm}{\otSemi\fontsize{9}{11}\selectfont\color{matAccent}\otCaps[8]{Partie 7 · Longueurs et masses}}\par\vspace{1mm}
\guideEntree{act-ms-39}{39}{Courts, moyens, longs}{Classer des objets selon leur longueur en trois groupes, quand les longueurs sont très différentes.}{Ranger des bandes de papier ou des bâtons en trois tas : les courts, les moyens, les longs.}{« Ce ruban est long, celui-ci est court, celui-là est entre les deux. »}{Trois rubans jaunes (courts), trois verts (moyens), trois rouges (longs).}{ms-39}
\guideEntree{act-ms-40}{40}{Du plus court au plus long}{Ranger des objets rectilignes du plus court au plus long et verbaliser le résultat.}{Ranger quatre bandes de papier de couleurs différentes du plus court au plus long, en les posant contre une règle.}{« Le vert est plus court que le bleu. Le rouge est le plus long. »}{Un point pour le vert, deux pour le bleu, trois pour le jaune, quatre pour le rouge.}{ms-40}
\guideEntree{act-ms-41}{41}{De même longueur}{Comparer des longueurs et reconnaître deux objets de même longueur.}{Comparer deux ficelles en les plaçant côte à côte, bout contre bout, au même point de départ.}{« Ce ruban est de même longueur que celui-là. Le bleu est plus long. »}{Premier cadre : le rouge et le vert (le bleu est plus long). Deuxième cadre : le jaune et l’orange (le violet est plus court).}{ms-41}
\guideEntree{act-ms-42}{42}{Le plus lourd}{Comparer les masses de deux objets très différents ; comprendre qu’un gros objet n’est pas forcément lourd.}{Soupeser une balle de tennis et une boule de pétanque, un sachet de coton et un sachet de sable de même taille.}{« Le ballon est plus gros, mais il est plus léger que la boule. La boule est plus lourde. »}{La boule (plus lourde que la plume), la boule (plus lourde que le grand ballon), la pastèque (plus lourde que la cerise).}{ms-42}
\guideEntree{act-ms-43}{43}{La balance}{Utiliser une balance de Roberval pour comparer deux masses.}{Poser deux objets sur une vraie balance de Roberval et observer le plateau qui descend.}{« Ce plateau est descendu : la pastèque est plus lourde que la pomme. »}{La pastèque, la boule, l’ours et la valise : à chaque fois, l’objet du plateau qui est descendu.}{ms-43}
\par\needspace{62mm}\vspace{6mm}{\otSemi\fontsize{9}{11}\selectfont\color{matAccent}\otCaps[8]{Partie 8 · Les motifs}}\par\vspace{1mm}
\guideEntree{act-ms-44}{44}{Continue la rangée}{Prolonger un motif répétitif et dire la règle utilisée.}{Faire des rangées avec des objets de la classe (crayon, gomme, crayon, gomme…) ; l’enfant continue et explique.}{« Étoile, rond, rond, étoile, rond, rond : après, il y a une étoile. »}{Étoile, rond, étoile ; étoile, rond, rond ; carré, carré, triangle ; rond, triangle, étoile. D’autres réponses sont possibles si l’enfant explique sa règle.}{ms-44}
\guideEntree{act-ms-45}{45}{L’intrus de la rangée}{Reconnaître un motif répétitif à ses régularités et trouver l’élément qui ne respecte pas la règle.}{Faire un collier qui alterne deux couleurs en glissant une erreur ; l’enfant la trouve et la corrige.}{« Étoile, rond, étoile, rond… Là, il y a deux ronds de suite : c’est l’intrus. »}{Première rangée : la sixième forme (une étoile à la place d’un rond). Deuxième : la huitième (un carré à la place d’un triangle). Troisième : la sixième (un rond à la place d’une étoile).}{ms-45}
\guideEntree{act-ms-46}{46}{Ça se répète ?}{Identifier parmi plusieurs rangées celles qui contiennent un motif répétitif.}{Écouter des rythmes frappés : certains se répètent (tap, tap, pause…), d’autres non.}{« Rond, carré, rond, carré : ça se répète. Dans l’autre, il n’y a pas de règle. »}{La première rangée (rond, carré) et la troisième (étoile, étoile, triangle).}{ms-46}
\guideEntree{act-ms-47}{47}{La tour de mémoire}{Mémoriser un motif répétitif simple pour le reproduire.}{Montrer une tour de briques (une bleue, deux jaunes…), la cacher, et la faire reconstruire de mémoire.}{« Une bleue, deux jaunes, une bleue, deux jaunes. »}{De bas en haut : bleu, jaune, jaune, bleu, jaune, jaune. Retenir « une bleue, deux jaunes » aide à mémoriser.}{ms-47}
\guideEntree{act-ms-48}{48}{Sous la carte}{Compléter un motif répétitif en trouvant l’élément manquant.}{Retourner une carte dans une rangée qui se répète et faire deviner ce qu’elle cache.}{« Étoile, rond, rond, étoile, rond, rond, étoile : sous la carte, il y a une étoile. »}{Une étoile, puis un triangle, puis un rond.}{ms-48}
\par\needspace{62mm}\vspace{6mm}{\otSemi\fontsize{9}{11}\selectfont\color{matAccent}\otCaps[8]{Partie 9 · Se repérer}}\par\vspace{1mm}
\guideEntree{act-ms-49}{49}{Entre les deux}{Situer des objets les uns par rapport aux autres : « entre ».}{Jeu de la file : « Mets-toi entre Léo et Sarah. » Ranger des jouets sur une étagère en suivant des consignes.}{« Le ballon est entre le chat et le lapin. »}{En haut, le ballon rouge ; en bas, le canard.}{ms-49}
\guideEntree{act-ms-50}{50}{En haut, au milieu, en bas}{Utiliser les mots « en haut », « en bas », « au milieu » pour situer des objets.}{Ranger des jouets sur une étagère et faire dire où se trouve chacun.}{« L’ours est tout en haut. La voiture est en bas. Le bateau est au milieu. »}{L’ours est entouré, la voiture barrée, le bateau colorié.}{ms-50}
\guideEntree{act-ms-51}{51}{Dans l’enclos}{Comprendre et utiliser « à l’intérieur » et « à l’extérieur ».}{Tracer un grand cercle à la craie dans la cour : « Entre dans le cercle, sors du cercle. »}{« Ce lapin est à l’intérieur de l’enclos. Ce canard est à l’extérieur. »}{Les quatre animaux à l’intérieur de la barrière. Les quatre autres sont à l’extérieur.}{ms-51}
\guideEntree{act-ms-52}{52}{Le labyrinthe}{Suivre un déplacement sur une représentation et le tracer.}{Faire un parcours dans la cour entre des plots, puis le refaire avec le doigt sur un dessin.}{« Je vais tout droit, je descends, je tourne… Ce chemin est fermé : je reviens. »}{Un seul chemin mène au fromage ; il est tracé en vert sur l’image du guide.}{ms-52}
\guideEntree{act-ms-53}{53}{Le chemin du chat}{Raconter un déplacement à partir d’une représentation, avec des repères et les mots de l’espace.}{Après un parcours en salle de motricité, le raconter en nommant les repères : « je passe près du banc, je tourne… ».}{« Le chat va tout droit, il descend près du champignon, il tourne, il passe à côté du sapin, il traverse le pont et il monte jusqu’à sa maison. »}{Le trait suit la route grise. Le récit doit citer les repères dans l’ordre (champignon, sapin, pont).}{ms-53}
